% Generated draft: requires mathematical and source review.
\documentclass[11pt]{article}
\usepackage[margin=25mm]{geometry}
\usepackage{fontspec}
\setmainfont{Latin Modern Roman}
\usepackage{amsmath,amssymb,mathtools,unicode-math}
\setmathfont{Latin Modern Math}
\usepackage{enumitem,needspace,xurl,hyperref}
\setlength{\parindent}{0pt}
\setlength{\parskip}{5pt}
\setlength{\emergencystretch}{3em}
\newcommand{\N}{\mathbb{N}}
\newcommand{\Z}{\mathbb{Z}}
\newcommand{\Q}{\mathbb{Q}}
\newcommand{\R}{\mathbb{R}}
\newcommand{\C}{\mathbb{C}}
\newcommand{\F}{\mathbb{F}}
\newcommand{\A}{\mathbb{A}}
\newcommand{\PP}{\mathbb{P}}
\newcommand{\E}{\mathbb{E}}
\newcommand{\eps}{\varepsilon}
\newcommand{\dd}{\,\mathrm{d}}
\begin{document}
\section*{UnsolvedMath ID: 30001216}
\textbf{Source problem code:} catalogue entry
\subsection{Definitions and mathematical statement}
Let \(M\) denote the localized equivariant cohomology representation associated in the supplied record with moduli spaces of parabolic sheaves on \(\PP^1\times\PP^1\). Write \(Y\) for the affine Yangian acting on \(M\), and write \(V_{\mathrm{univ}}\) for the intended universal Verma module over \(U(\widehat{\mathfrak{gl}}_n)\), the universal enveloping algebra of the indicated affine Lie algebra. Here these symbols retain the source's intended meanings: the record does not specify the allowed ranks \(n\), the precise affine Lie algebra convention, the moduli spaces, the equivariance group, the localization, the Yangian parameters, or the coefficient ring and highest weight defining \(V_{\mathrm{univ}}\).

The proposed identification is
\[
M\ \cong\ V_{\mathrm{univ}}
\qquad\text{as affine Yangian modules},
\]
with the additional requirement that the affine Gelfand--Tsetlin fixed-point basis specialize to bases of integrable modules.

This is an unresolved formulation of the source question, conditional on specifying the omitted data. In particular, a \(U(\widehat{\mathfrak{gl}}_n)\)-module structure alone does not specify a \(Y\)-module structure on \(V_{\mathrm{univ}}\); the comparison requires that additional structure and compatible coefficient rings. The basis requirement also needs the intended integrable modules, specialization parameters and maps, and a prescription for passing from localized fixed-point classes to specialized basis vectors. The record does not state whether every integrable module in some class is intended, or which class that would be. Consequently, neither an unconditional module-isomorphism assertion nor a quantified specialization assertion is fully determined by the supplied record.
\subsection{Short English statement}
The question asks whether the localized cohomology representation of parabolic-sheaf moduli spaces on \(\PP^1\times\PP^1\) realizes the intended universal affine Verma module, compatibly with affine Yangian actions, and whether its fixed-point basis specializes to bases of the intended integrable modules. The actions, parameters, and specialization procedure need further specification.
\subsection{Sources}
\begin{itemize}
\item Supplied local UnsolvedMath registry record, permanent ID \(30001216\), source problem code \texttt{catalogue entry}, title \emph{Yangian Realization of Universal Affine Verma Modules}, in the Oberwolfach Reports Open Problems collection. The record gives proposed year 2009 and names no individual proposer.
\item \emph{Enveloping Algebras and Geometric Representation Theory} (2009), Oberwolfach Reports, pp.\ 801--866, DOI 10.4171/owr/2009/15: \url{https://doi.org/10.4171/owr/2009/15}. The supplied background describes these reports as documenting mathematical workshops and their open-problem discussions.
\item Supplied literature reference labeled ``Feigin--Finkelberg--Negut--Rybnikov primary paper'': \url{https://doi.org/10.1007/s00029-011-0059-x}.
\item Supplied literature reference labeled ``primary preprint'': \url{https://arxiv.org/abs/0812.4656}.
\item Supplied literature reference labeled ``original OWR report'': \url{https://doi.org/10.4171/owr/2009/15}.
\item Literature assessment dated exactly \texttt{2026-08-21}: the supplied curation labels the problem partially solved, reports the affine Yangian action, fixed-point formulas, and affine Gelfand--Tsetlin algebra surjection as established, and describes the stronger identification of the full localized cohomology module with the universal affine Verma module as conjectural in the primary report. This is supplied dated curation, not an independently verified assessment; its time-sensitive status and cited claims require independent verification.
\end{itemize}
\end{document}
